\Askhsh{34438_SOLUTION}{1}
\begin{ylh}{1}
    \item [Διαφορικός Λογισμός]
\end{ylh}

\OnPar{ΛΥΣΗ}

\begin{alist}
\item α) Η συνάρτηση $f$ είναι δύο φορές παραγωγίσιμη στο $\mathbb{R}$ με:
\[ f'(x) = (2x^3 - 15x^2 + 24x)' = 6x^2 - 30x + 24 \]
και
\[ f''(x) = (6x^2 - 30x + 24)' = 12x - 30. \]
Επίσης $f'(x) = 0 \Leftrightarrow 6(x^2 - 5x + 4) = 0 \Leftrightarrow x^2 - 5x + 4 = 0 \Leftrightarrow x = 1 \text{ ή } x = 4$.
και $f''(x) = 0 \Leftrightarrow 12x - 30 = 0 \Leftrightarrow x = \frac{30}{12} = \frac{5}{2}$.

\item β) Λόγω του ερωτήματος (α) είναι: $f'(x) > 0 \Leftrightarrow x < 1 \text{ ή } x > 4$ οπότε έχουμε τον παρακάτω πίνακα μεταβολών:
\[
\begin{array}{|c|c|c|c|c|c|}
\hline
x & -\infty & 1 & 4 & +\infty \\
\hline
f'(x) & + & 0 & - & 0 & + \\
\hline
f(x) & \nearrow & \text{TM} & \searrow & \text{TE} & \nearrow \\
\hline
\end{array}
\]
Επομένως η $f$ είναι γνησίως αύξουσα στα διαστήματα $(-\infty, 1]$, $[4, +\infty)$ και γνησίως φθίνουσα στο διάστημα $[1, 4]$.
Στο $x_1 = 1$ εμφανίζει τοπικό μέγιστο το $f(x_1) = f(1) = 11$ και στο $x_2 = 4$ εμφανίζει τοπικό ελάχιστο το $f(x_2) = f(4) = -16$.

\item γ) Λόγω του ερωτήματος (α) είναι: $f''(x) > 0 \Leftrightarrow x > \frac{5}{2}$ οπότε έχουμε τον παρακάτω πίνακα μεταβολών:
\[
\begin{array}{|c|c|c|c|}
\hline
x & -\infty & \frac{5}{2} & +\infty \\
\hline
f''(x) & - & 0 & + \\
\hline
f(x) & \curvearrowleft & \text{Σ.Κ.} & \curvearrowright \\
\hline
\end{array}
\]
Επομένως η $f$ είναι κοίλη στο διάστημα $(-\infty, \frac{5}{2}]$ και κυρτή στο διάστημα $[\frac{5}{2}, +\infty)$.
Επίσης στο $x_3 = \frac{5}{2}$ η $f$ εμφανίζει σημείο καμπής.
\end{alist}